\section{Discussion}
\label{sec:disc}

\emph{Dispersion and line width.} The reciprocal dispersion was found to be \DispRatio{} times the design value, and the lines (\FWHMKr~px) are narrower than the expected slit image (\SlitImgPx~px). Both effects can be explained by a prism rotated away from minimum deviation: at $\theta_1\approx\ThetaOff^\circ$ the anamorphic factor $r=\AnamOff$ reproduces the measured line width and accounts for a factor \DispRatioOff{} in $P$. The remaining $\sim$20\% is attributed to the real focal length of the camera doublet and the detector position, which were not measured.

\emph{Resolution.} $R$ lies above the design curve because the narrower lines more than compensate the lower dispersion. The instrument is slit limited ($R_{dif}$ is ten times larger) and the $\sim$3~px FWHM satisfies Nyquist sampling. The Ar values are the least reliable, due to blends and a less uniform illumination of the slit.

\emph{Range.} The flash range is limited by the source; the instrument covers at least \RangeDetMin--\RangeLastLine~nm. In the near infrared $P>2$~nm/px (700--970~nm span only \NPxNIR~px), so the useful range is the visible.

\emph{Error sources.} The blue calibration is an extrapolation (no lamp lines below 557~nm) checked only with LEDs of unknown wavelength, and the lamps were identified by their line pattern. The 8-bit format, the missing darks and the saturated flash affect the profiles more than the centroids.
